%\usepackage[margin=0.5cm]{geometry}
\usepackage[left=1cm,right=1cm,top=1cm,bottom=2cm]{geometry}
% Note: inputenc not needed with lualatex (handles UTF-8 natively)
%\usepackage[light, largesmallcaps]{kpfonts}
\usepackage{mwe}
\usepackage{multicol}
\usepackage{pgfkeys}

%\usepackage{geometry}
\usepackage{lipsum}
\usepackage{microtype}
\usepackage{grid-system}
\usepackage{babel}
\usepackage{listings}
\usepackage{graphicx}
\usepackage{import}
\usepackage{verse}
\usepackage{setspace}

\usepackage{aurical}
% \usepackage[T1]{fontenc}
%\usepackage{css-colors}
\usepackage{latexcolors}
% xcolor is loaded via tikz in main.tex with svgnames option
% \usepackage{tabularray}
\usepackage{nicematrix}

% we want only the chapters in the toc, for books. No toc for rsongs
%\setcounter{tocdepth}{1}
\usepackage{amssymb}
\usepackage[default]{frcursive}

\usepackage{fancyhdr}
\usepackage{lastpage}
\pagestyle{fancy}
\fancyhf{}
%\usepackage{tabularx}
\usepackage[most]{tcolorbox}

\usepackage{booktabs}
\usepackage{array}
\usepackage{paracol}

\usepackage{hyperref}
\hypersetup{
	colorlinks,
	citecolor=black,
	filecolor=black,
	linkcolor=black,
	urlcolor=black,
	bookmarks=false,
	hypertexnames=false
}

\usepackage{units}

\usepackage[most]{tcolorbox}
\usepackage{tikz, adjustbox}

%\setlength{\headheight}{10pt}
\usepackage[normalem]{ulem}
\renewcommand{\headrulewidth}{0pt}
\renewcommand{\footrulewidth}{2pt}

%\lhead{\includegraphics[width=1cm]{example-image-a}}
\rhead{}



\rfoot{\thepage/\pageref{LastPage}}
%\cfoot{\today}
%\cfoot{derni\`ere modif le \songlastupdate, g\'en\'er\'e le \songtoday }
%\cfoot{derni\`ere modif le \songlastupdate }

\definecolor{bg}{HTML}{ADFF2F}
%\definecolor{bgred}{\color{Seashell3}}

% Song metadata from YAML (handlebars template)
\newcommand{\songtitle}{ {{{song.info.title}}} }
\newcommand{\songauthor}{ {{{song.info.author}}} }
\newcommand{\songtempo}{ {{{song.info.tempo}}} }
\newcommand{\songlastupdate}{ \today }

\usepackage{fontspec}
\setmainfont{Garamond Libre}

% pour les lignes pointillées dans les tableaux
% \usepackage{arydshln}


% #3 is always \input{lyrics/<section>}, and that file is empty for a section
% whose lyrics are not written yet. Typeset it into a scratch box first: when it
% comes out empty, skip both the line break and the 16pt block, so a run of
% lyric-less sections stacks its colour bars tightly instead of leaving a blank
% line under each one. The box is measured only, then thrown away and #3 is
% typeset for real.
\newsavebox{\basecoupletprobe}

% Name of the chord-chart coordinate at the center of bar #1, for add.tikz:
% \path (\BAR{4}) ... The argument is an integer expression, so \BAR{4+1}
% works too. song.tikz defines one BAR_<n> coordinate per bar.
\newcommand*{\BAR}[1]{BAR_\the\numexpr#1\relax}

% Name of the chord-chart coordinate of section #1 (its id in song.yml), for
% add.tikz: \path (\SECTION{refrain1}) ... It sits on the right edge of the
% section's bar grid, halfway down the section.
\newcommand*{\SECTION}[1]{SECTION_#1}

\newcommand{\basecouplet}[3]{
	\vspace{6pt}
	\fcolorbox{gray}{#1}{\parbox{\dimexpr\columnwidth-2\fboxsep-2\fboxrule}{
		\fontsize{18pt}{18pt}\selectfont
		#2
	}}%
	\setbox\basecoupletprobe=\vbox{\hsize=\columnwidth
		\fontsize{16pt}{16pt}\selectfont
		#3
	}%
	\ifdim\ht\basecoupletprobe>0pt
		\\[6pt]
		{
			{
					\fontsize{16pt}{16pt}\selectfont
					#3
				}
		}
	\fi
}


\newcommand{\xxcfoot}[2]{
\cfoot\{{
		\Fontskrivan\bfseries\slshape
		\fontsize{14pt}{14pt}\selectfont
		\color{blue}
		#2
	}
	{
		\cursive\bfseries
		\fontsize{10pt}{10pt}\selectfont
		\color{orange}
		#1
	}
{{#if song.meta.date}}
; {\fontsize{10pt}{10pt}\selectfont derni\`ere modif le \songlastupdate}
{{/if}}
}
}


\newcommand{\makesongtitle}{}

\newcommand{\xxmakesongtitle}[3] {
	\begin{center}
		{
			\Fontskrivan\bfseries\slshape
			\fontsize{60pt}{50pt}\selectfont
			\color{blue}
			#1
		} \\

		{
		\textcursive{
			\bfseries%
			\fontsize{20pt}{10pt}\selectfont%
			\color{orange}%
			( #2 )%
		}
		}

		{
		%\Fontskrivan\bfseries\slshape
		\fontsize{10pt}{5pt}\selectfont
		\color{blue}
		#3 BPM
		}\\
	\end{center}

}





\newcommand{\songlyrics}{}


\newcommand{\songly}[1]{ %
	\subimport{#1.output}{#1}
}

\newcommand{\songbookcomment}[1]{%
	#1
}

\newcommand{\songword}[1]{%
	%\tikz[]\node [direct] (N2) {#1};%
	#1%
}

% songworkd first beat
\newcommand{\songwordfb}[1]{%
	%\tikz[]\node [firstbar] (N1) {#1};%
	\fcolorbox{black}{red!20}{#1}%
}

%levee
\newcommand{\songwordl}[1]{%
	%\fcolorbox{black}{blue!10}{$[$\it #1$]$}%
	%\tikz[]\node [levee] (N1) {#1};%
	{\color{blue!60}{$[${\it #1}$]$}}%
}

\newcommand{\songwordcount}[1]{
	\textsubscript{\tiny\color{red}#1}%
}

